\section{Introduction}

Distributed Asynchronous Object Storage (DAOS) is a high-performance, open-source parallel I/O stack built on Non-Volatile Memory Express (NVMe) and Remote Direct Memory Access (RDMA) technologies~\cite{daos-exascale-storage2016sc, daos-intro2020, daos-dram-intro2023, daos-github}. It replaces conventional POSIX file and directory abstractions with a versioned, key-value object model and provides direct object and filesystem interfaces for high-performance applications. Since its initial deployment in 2020, prior studies have demonstrated strong DAOS scalability and performance, showing that it can match or outperform traditional parallel file systems across a variety of platforms and workloads~\cite{daos-vs-lustre2022pdsw, daos-arm2023SCAsia, daos-scalability2023SCAsia, daos-perf-explore2024scw, DAOS-aurora-20-node-test2024scw, XianheSun2023evaluation}. DAOS has also been studied for reducing I/O contention and accelerating data access in weather-prediction workflows, scientific data formats, and machine-learning pipelines~\cite{hdf5-daos2022tpds, daos-adios2024ISC, hpc-weather2023ipdps, manubens2024reducing, daos-smartnic2025scw, devarajan2021dlio, mlperf-storage-github, brim2025dug-daos-aio}. This work is motivated by using DAOS as a burst buffer for scientific streaming systems, where data from many scientific instruments arrive directly at high-performance computing (HPC) facilities and must be absorbed continuously at line rate without intermediate staging~\cite{ejfat2022indis}.

Scientific streaming places a different set of I/O demands on storage than the checkpoint- and training-oriented workloads emphasized in much of the prior DAOS literature. First, many instrument streams produce data concurrently, requiring high throughput under intra-node I/O concurrency and as the number of client nodes increases. Second, continuous data production makes latency important in addition to aggregate throughput. An individual slow write can delay a producer and introduce backpressure into the streaming pipeline, so mean latency alone does not capture the relevant behavior. Third, a streaming workflow spans different access phases, including write-dominated data production, mixed read/write online analysis, and read-dominated post-processing. Finally, raw detector events can be relatively small, approximately 1~MiB each in our target use case. The resulting performance therefore depends not only on aggregate bandwidth, but also on I/O request size, concurrency, access order, and how small events are organized into storage objects.


Existing DAOS evaluations establish its high peak throughput and scalability, but these results do not fully characterize how DAOS behaves under the combination of throughput, latency, and data-access requirements above. We therefore use controlled \textit{fio} and IOR experiments to isolate the storage parameters that are most relevant to scientific streaming. Rather than evaluating an end-to-end streaming application, our goal is to identify how these parameters affect DAOS performance and to derive configuration guidance for future streaming deployments. We characterize two different DAOS deployments, including a large-scale deployment on Aurora~\cite{aurora-intro2025}, and address the following research questions. \textbf{RQ1:} How do I/O request size, concurrency, DAOS access method, and client-node scale affect DAOS throughput? (Sec~\ref{sec:rq1_throughput}) \textbf{RQ2:} At fixed total I/O concurrency, how does the process/queue-depth configuration affect tail latency? (Sec~\ref{sec:rq2_latency}) \textbf{RQ3:} How do read/write mix, access order, and object size and count affect DAOS data access performance? (Sec~\ref{sec:rq3_us}) %We additionally compare DAOS with NFS to provide a conventional network file-system reference for request-size sensitivity. 

Our evaluation yields the following main findings:


\begin{comment}
    
\begin{itemize}
\item I/O request size relative to the DAOS chunk size is the first-order performance factor. Chunk-aligned requests saturate the NIC bandwidth with modest concurrency, whereas smaller requests remain IOPS- and CPU-bound.

\item At fixed total concurrency, the process/queue-depth split leaves throughput unchanged but can inflate P99.99 latency to 26$\times$ the median. More processes with shallower queues minimize tail latency.

\item DAOS scales near-linearly to multi-TiB/s across client nodes for large objects, but per-object overhead causes many-small-object reads to plateau at $3$--$5\times$ lower, motivating event aggregation before ingestion.
\end{itemize}


The remainder of this paper is organized as follows. Section~\ref{sec:background} reviews the DAOS architecture and prior performance studies. Section~\ref{sec:methods} describes our experimental platforms and benchmark configurations. Section~\ref{sec:results} presents our measurement results and findings. Section~\ref{sec:conclusion} concludes the paper and discusses future work.

\end{comment}






\begin{itemize}
\item I/O request size is the primary throughput factor under intra-node concurrency. MiB-scale requests approach client NIC bandwidth with modest concurrency, while small requests remain CPU- and IOPS-bound. On Aurora, DAOS scales to multi-TiB/s across client nodes.
\item At fixed total I/O concurrency and comparable throughput, the process/queue-depth configuration substantially changes P99.99 completion latency. Deep queues increase completion tail latency, while very shallow queues increase worst-case submission latency.
\item Sequential and random throughput largely converge at large request sizes, while small requests remain sensitive to access order. At scale, reading the same data volume from many small objects achieves 3--5$\times$ lower throughput than reading from large objects.
\end{itemize}